\begin{center}
Problema G

Fabiola e os cafés

{\tiny Por Cleon Xavier Pereira Júnior}

Timelimit: $1$	
\end{center}			

{\footnotesize
Professora Fabiola é fascinada por cafés de qualidade. Em 2022, a professora Fabiola pegou seu Ford Ka e, com sua filha, foi aventurar-se pelo sul de minas para experimentar os melhores cafés do país. Em cada produtora de cafés especiais que Fabiola visitava, ela levava um pacote. Os pacotes tinham diferentes valores e pesos. 

Acontece que, na empolgação, a professora Fabiola esqueceu que tinha um Ford Ka, uma filha e umas bagagens. Já sabem o que aconteceu né? Eis que a professora Fabiola tinha um problema de excesso de carga. Ao término da viagem, quando ela colocou tudo no carro, a professora se deparou com um carro rebaixado, mas não era por estilização, era por questão de peso mesmo. 

Observando a inviabilidade da situação, a professora Fabiola notou que deveria deixar uns cafés para trás. Claro que ela gostaria de levar o máximo de cafés possíveis e que tivesse o menor prejuízo. Para isso, ela estava disposta, se necessário, a abrir alguns pacotes de café e levar só parte desse, para ter a possibilidade de levar a capacidade máxima.

No momento, a professora Fabiola só está preocupada em saber o tamanho do prejuízo financeiro que ela terá na melhor situação. Ou seja, quanto ela gastou e não conseguirá levar por excesso de peso. Vamos ajudar a professora fazendo este cálculo para ela?


\vspace{0.3cm}
\textbf{Entrada}
Cada entrada contém:
\begin{itemize}
    \item um inteiro $X$ que corresponderá aos $M$ pacotes de cafés;
    \item um número $Y$, correspondente à capacidade de carga disponível no carro para os cafés;
    \item $M$ linhas com dois números, separados por um espaço, sendo o peso, seguido do valor.
\end{itemize}

\vspace{0.3cm}
\textbf{Saída}
O algoritmo deve retornar um valor numérico, com duas casas decimais, correspondendo ao prejuízo financeiro que a professora terá.
\begin{table}[h]
    \centering
    \begin{tabular}{|p{8cm}|p{7cm}|}
        \hline
        \begin{center}
            \vspace{-0.3cm}
            \textbf{Exemplos de Entrada}
            \vspace{-0.3cm}
        \end{center} 
         &  
        \begin{center}
            \vspace{-0.3cm}
            \textbf{Exemplos de Saída}
            \vspace{-0.3cm}
        \end{center} \\ \hline
         \texttt{5} & \texttt{1200.00} \\  
         \texttt{50} & \texttt{} \\  
         \texttt{40 840} & \texttt{} \\
         \texttt{30 600} & \texttt{} \\
         \texttt{20 400} & \texttt{} \\
         \texttt{10 100} & \texttt{} \\
         \texttt{20 300} & \texttt{} \\
         \hline
                  \texttt{3} & \texttt{2653.12} \\  
         \texttt{450} & \texttt{} \\  
         \texttt{420 1500} & \texttt{} \\
         \texttt{450 1200} & \texttt{} \\
         \texttt{480 1550} & \texttt{} \\
\hline
    \end{tabular}
    \caption{Exemplos de entradas e saídas}
    \label{tab:inputProblem}
\end{table}}

\newpage

\begin{center}
\noindent
\lstinputlisting{abobora_23/G/G3.cpp}
\end{center}